\subsection{Locally trivial fiber spaces with base B × I}

PROPOSITION 2. --- Let B be a paracompact topological space and let (E, p) be a locally trivial fiber space over B × I. Set E₁ = \(\overline{p}^{-1}(B \times \{1\})\) and denote by \(p_{1}: E_{1} \to B\) the map \(pr_{1} \circ p | E_{1}\). Then the B × I-spaces (E, p) and (E₁ × I, \(p_{1} \times Id_{I}\)) are isomorphic.

Let us first prove two lemmas.

Lemma 2. --- Let \(\alpha\), \(\beta\), \(\gamma\) be real numbers such that \(\alpha \leqslant \beta \leqslant \gamma\); let B be a topological space and let \(p: \mathrm{E} \to \mathrm{B} \times [\alpha, \gamma]\) be a continuous map. Set \(\mathrm{B}_0 = \mathrm{B} \times [\alpha, \beta]\), \(\mathrm{B}_1 = \mathrm{B} \times [\beta, \gamma]\), \(\mathrm{E}_0 = p^{-1}(\mathrm{B}_0)\), \(\mathrm{E}_1 = p^{-1}(\mathrm{B}_1)\) and denote by \(p_0: \mathrm{E}_0 \to \mathrm{B}_0\), \(p_1: \mathrm{E}_1 \to \mathrm{B}_1\) the maps induced by \(p\). If \((\mathrm{E}_0, p_0)\) and \((\mathrm{E}_1, p_1)\) are trivializable fiber spaces, the same holds for \((\mathrm{E}, p)\).

Let \(g_0 \colon \mathrm{E}_0 \to \mathrm{B}_0 \times \mathrm{F}_0\) and \(g_1 \colon \mathrm{E}_1 \to \mathrm{B}_1 \times \mathrm{F}_1\) be trivializations of \(\mathrm{E}_0\) and \(\mathrm{E}_1\), respectively. Let us denote by \(g_0'\) and \(g_1'\) the trivializations of the \(\mathrm{B} \times \{\beta\}\)-fiber space \(p^{-1}(B \times \{\beta\})\) deduced from \(g_0\) and \(g_1\) by restriction. The map \(h = g_0' \circ (g_1')^{-1}\) is a \(B \times \{\beta\}\)-isomorphism of \(B \times \{\beta\} \times F_1\) onto \(B \times \{\beta\} \times F_0\). We define a continuous map \(h'\) from \(B \times F_1\) to \(F_0\) by setting \(h'(a, y) = \mathrm{pr}_3 \circ h(a, \beta, y)\) for \((a, y) \in B \times F_1\). For \((a, t, y) \in B \times [\beta, \gamma] \times F_1\), set \(H(a, t, y) = (a, t, h'(a, y))\). The map H thus defined is a \((B \times [\beta, \gamma])\)-isomorphism of \(B \times [\beta, \gamma] \times F_1\) onto \(B \times [\beta, \gamma] \times F_0\), and one has \(g_0 | p^{-1}(B \times \{\beta\}) = H \circ g_1 | p^{-1}(B \times \{\beta\})\). There therefore exists a continuous map \(g \colon E \to B \times [\alpha, \gamma] \times F_0\) such that \(g | E_0 = g_0\) and \(g | E_1 = H \circ g_1\). The map \(g\) is an isomorphism of \(B \times [\alpha, \gamma]\)-spaces, hence E is trivializable.

Lemma 3. --- Let B be a topological space and let (E, p) be a locally trivial fiber space over B×I. Every point a of B has a neighborhood V such that the V × I-space \(E_{V \times I}\) is trivializable.

Let \(a\) be a point of B; for every point \(t\) of \(\mathbf{I}\), there exists an open neighborhood \(\mathrm{W}_t\) of \(t\) in \(\mathbf{I}\) and a neighborhood \(\mathrm{V}_t\) of \(a\) in B such that E is trivializable over \(\mathrm{V}_t \times \mathrm{W}_t\). There then exists an integer \(n > 0\) and, for every integer \(i\) such that \(1 \leqslant i \leqslant n\), a point \(t_i\) of \(\mathbf{I}\) such that the interval \([\frac{i-1}{n}, \frac{i}{n}]\) is contained in \(\mathrm{W}_{t_i}\) (III, p. 272, lemma 4). Set \(\mathrm{V} = \cap_{1 \leqslant i \leqslant n} \mathrm{V}_{t_i}\). The fibered space E is trivializable over \(\mathrm{V} \times [\frac{i-1}{n}, \frac{i}{n}]\) for every integer \(i\) such that \(1 \leqslant i \leqslant n\). Lemma 3 then follows from Lemma 2 by induction on \(n\).

Now let us prove the proposition. By Lemma 3, there exists an open covering \((\mathrm{U}_j)_{j\in \mathrm{J}}\) of B such that, for every \(j\in \mathrm{J}\), E is trivializable over \(\mathrm{U}_j\times \mathbf{I}\). Since the space B is paracompact, we may assume the covering \((\mathrm{U}_j)_{j\in \mathrm{J}}\) locally finite (TG, IX, p. 49) and choose a covering \((\mathrm{A}_j)_{j\in \mathrm{J}}\) of B where, for every \(j\in \mathrm{J}\), the set \(\mathrm{A}_j\) is closed in B and contained in \(\mathrm{U}_j\) (TG, IX, p. 49, prop. 4 and p. 48, cor. 1).

For every open subset U of B, denote by \(\mathcal{F}(\mathrm{U})\) the set of \(\mathrm{U}\times\mathbf{I}\) -isomorphisms of \(\stackrel{-1}{p}(\mathrm{U}\times\mathbf{I})\) onto \(\stackrel{-1}{p_{1}}(\mathrm{U})\times\mathbf{I}\) which induce the identity map of \(\stackrel{-1}{p}(\mathrm{U}\times\{1\})\) onto \(\stackrel{-1}{p_{1}}(\mathrm{U})\times\{1\}\) . For every pair (V, U)

of open sets of B such that U ⊂ V, denote by \(r_{\mathrm{UV}} \colon \mathcal{F}(\mathrm{V}) \to \mathcal{F}(\mathrm{U})\) the map which, to a \((\mathrm{V} \times \mathbf{I})\) -isomorphism \(g \colon \stackrel{-1}{p}(\mathrm{V} \times \mathbf{I}) \to \stackrel{-1}{p_0}(\mathrm{V}) \times \mathbf{I}\) , associates the \((\mathrm{U} \times \mathbf{I})\) -isomorphism deduced from g by passing to subspaces. The pair \(\mathcal{F} = ((\mathcal{F}(\mathrm{U})), (r_{\mathrm{UV}}))\) is a sheaf on B (I, p. 45, example 4). To prove the proposition, it suffices to prove that the sheaf F is soft (I, p. 64).

Let \(j \in \mathrm{J}\) , let A be a closed subset of \(\mathrm{A}_j\) , let V be an open set in B such that \(\mathrm{A} \subset \mathrm{V} \subset \mathrm{U}_j\) and let \(g\) be an element of \(\mathcal{F}(\mathrm{V})\) . There exists an open neighborhood W of A such that \(\overline{\mathrm{W}} \subset \mathrm{V}\) , because a paracompact space is normal (TG, IX, p. 49, prop. 4). We shall prove that there exists an element \(g'\) of \(\mathcal{F}(\mathrm{U}_j)\) such that \(g' | \mathrm{W} = g | \mathrm{W}\) . Corollary 2 (I, p. 65) of prop. 6 then implies that the sheaf \(\mathcal{F}\) is soft.

Since the B × I-space E is trivializable over \(U_{j} \times I\), we may assume that \(\overline{p}^{-1}(U_{j} \times I) = U_{j} \times I \times F\), where F is a topological space. The element g of \(\mathcal{F}(V)\) is then a V × I-isomorphism of V × I × F onto itself which induces the identity map of V × \{1\}×F. Apply cor. 5 of IV, p. 439, to the spaces \(X' = X\), \(X = U_{j}\), A = overline{W}, \(U = V\), and \(Y = Z = F\), to the map g: V 	imes I 	imes F 	o V 	imes I 	imes F, and to the identity map of \(U_{j} \times \{1\} \times F\). There therefore exists a \(U_{j} \times I\)-isomorphism \(g'\) of \(U_{j} \times I \times F\) onto itself which induces the identity map of \(U_{j} \times \{1\} \times F\) and which coincides with g on \(\overline{W} \times I \times F\), and a fortiori on W × I × F. Hence the proposition.

COROLLARY 1. --- Let B be a paracompact topological space and let (E, p) be a locally trivial fiber space over B × I (I, p. 71, corollary 2). For \(t \in I\), denote by \((E_{t}, p_{t})\) the locally trivial fibered B-space \(i_{t}^{*}E\), where \(i_{t} \colon B \to B \times \{t\}\) is the map \(b \mapsto (b, t)\). The locally trivial fibered B-spaces \(E_{0}\) and \(E_{t}\) are isomorphic for every \(t \in I\).

COROLLARY 2. --- Let B and B' be topological spaces, let E be a locally trivial fiber bundle over B. Let \(f_0\) and \(f_1\) be continuous maps from B' into B. Suppose that the space B' is paracompact. If the maps \(f_0\) and \(f_1\) are homotopic, the locally trivial fibered B'-spaces \(f_0^*E\) and \(f_1^*E\) are isomorphic.

Let \(\sigma\colon\mathrm{B}^{\prime}\times\mathbf{I}\to\mathrm{B}\) be a homotopy joining \(f_{0}\) to \(f_{1}\) and denote by \(\mathrm{E}^{\prime}\) the \(\mathrm{B}^{\prime}\times\mathbf{I}\)-space \(\sigma^{*}\mathrm{E}\); it is a locally trivial fibered space. Denote by \(i_{0}\) and \(i_{1}\) the maps from \(\mathrm{B}^{\prime}\) into \(\mathrm{B}^{\prime}\times\mathbf{I}\) given by \(x\mapsto(x,0)\) and \(x\mapsto(x,1)\). By Cor. 1, the \(\mathrm{B}^{\prime}\)-fibered spaces \(i_{0}^{*}\mathrm{E}^{\prime}\) and \(i_{1}^{*}\mathrm{E}^{\prime}\) are isomorphic. Since \(\sigma\circ i_{0}=f_{0}\), the \(\mathrm{B}^{\prime}\)-space \(i_{0}^{*}\mathrm{E}^{\prime}\) is identified with \(f_{0}^{*}\mathrm{E}\);

Similarly, the B'-space \(i_{1}^{*}E'\) is identified with \(f_{1}^{*}E\). Consequently, the B'-fibered spaces \(f_{0}^{*}E\) and \(f_{1}^{*}E\) are isomorphic, which is what was to be proved.

COROLLARY 3. --- Let B be a paracompact topological space. If B is homotopy equivalent to a point, then every locally trivial fiber bundle with base B is trivializable.

COROLLARY 4. --- Let B and B' be topological spaces, let (E, p) be a locally trivial fiber space over B, let f be a continuous map from B' into B, let \(\sigma: \mathrm{B}' \times \mathbf{I} \to \mathrm{B}\) be a homotopy with initial map f, and let \(\widetilde{f}\) be a continuous lift of f to E. If the space B' is paracompact, there exists a homotopy \(\widetilde{\sigma}: \mathrm{B}' \times \mathbf{I} \to \mathrm{E}\) with initial map \(\widetilde{f}\) that is a lift of \(\sigma\) to E.

Let \((\mathrm{E}^{\prime},p^{\prime})\) be the \(B^{\prime}\times I\)-space deduced from \((\mathrm{E},p)\) by base change along the map \(\sigma\colon B^{\prime}\times I\to B\). It is a locally trivial fiber space (I, p. 71, cor. 2), and the map \(s\colon B^{\prime}\times\{0\}\to E^{\prime}\) defined by \(s((a,0))=((a,0),\widetilde{f}(a))\) is continuous for \(a\in B^{\prime}\) and is a continuous section of \(p^{\prime}\) over \(B^{\prime}\times\{0\}\). By Prop. 2, there exists a continuous section \(\widetilde{s}\) of \(p^{\prime}\) which extends s. The map \(\widetilde{\sigma}=\operatorname{pr}_{2}\circ\widetilde{s}\) has the required property.

